\documentclass[11pt]{scrartcl}
\usepackage[utf8]{inputenc}
\usepackage[sexy]{evanchen} 
\usepackage{graphicx}
\graphicspath{{8 Outsource K Math/Gambar/}}


\title{Keterbagian}
\author{Azzam and Akeyla}
\date{Maret 2023}

\begin{document}

\maketitle
\section{Keterbagian}
Diberikan dua buah bilangan bulat $a$ dan $b$ dengan $a \neq 0$. Kita menyebut $a$ membagi $b$ apabila terdapat suatu bilangan bulat $c$ sedemikian hingga $b = ca$. Kita akan menotasikan ini dengan $a|b$.

\subsection{Sifat keterbagian terhadap bilangan-bilangan kecil}
\begin{enumerate}
    \item Sebuah bilangan asli habis dibagi 2 apabila digit terakhirnya habis dibagi oleh 2
    \item Sebuah bilangan asli habis dibagi 4 apabila dua digit terakhirnya habis dibagi oleh 4
    \item Sebuah bilangan asli habis dibagi 8 apabila tiga digit terakhirnya habis dibagi oleh 8
    \item Secara umum, sebuah bilangan asli habis dibagi oleh $2^n$ apabila $n$ digit terakhirnya habis dibagi oleh $2^n$ untuk setiap bilangan asli $n$
    \item Sebuah bilangan asli habis dibagi oleh 5 apabila digit terakhirnya habis dibagi oleh 5 (yakni 0 atau 5)
    \item Sebuah bilangan asli habis dibagi oleh 25 apabila dua digit terakhirnya habis dibagi oleh 25 (yakni 00,25,50 atau 75)
    \item Secara umum, sebuah bilangan asli habis dibagi oleh $5^n$ apabila $n$ digit terakhirnya habis dibagi oleh $5^n$ untuk setiap bilangan asli $n$
    \item Sebuah bilangan asli habis dibagi oleh $3$ apabila jumlah digit bilangan tersebut habis dibagi $3$
    \item Sebuah bilangan asli habis dibagi 9 apabila jumlah digit bilangan tersebut habis dibagi 9
    \item Sebuah bilangan asli habis dibagi 11 apabila selisih jumlah digit di tempat ganjil dan jumlah digit di tempat genap juga habis dibagi 11
    \item Sebuah bilangan habis dibagi $10^n$ apabila n digit terakhirnya sama dengan 0
    \item Sebuah bilangan habis dibagi $ab$ untuk bilangan asli $a$ dan $b$ yang relatif prima (tidak punya faktor bersama selain 1) apabila bilangan tersebut habis dibagi oleh $a$ dan $b$
    
    Untuk seterusnya, untuk sebuah bilangan asli $n$, notasikan $f(n)$ sebagai bilangan asli yang terbentuk dari menghapus digit terakhir dari $n$ jika $n$ memiliki minimal dua digit dan sama dengan $0$ jika $n$ hanya punya satu digit. Notasikan pula $g(n)$ sebagai digit terakhir dari $n$. Sebagai contoh $f(421)=42$ dan $g(421)=1$

    \item Sebuah bilangan asli $n$ habis dibagi oleh $7$ apabila $f(n)-2g(n)$ juga habis dibagi oleh $7$
    \item Sebuah bilangan asli $n$ habis dibagi oleh $13$ apabila $f(n)-9g(n)$ juga habis dibagi oleh $13$
    \item Sebuah bilangan asli $n$ habis dibagi oleh $17$ apabila $f(n)-5g(n)$ juga habis dibagi oleh $17$
    \item Sebuah bilangan asli $n$ habis dibagi oleh $19$ apabila $f(n)-17g(n)$ juga habis dibagi oleh $19$
    \item Sebuah bilangan asli $n$ habis dibagi oleh $23$ apabila $f(n)-16g(n)$ juga habis dibagi oleh $23$
\end{enumerate}

\subsection{Latihan Sifat Keterbagian Bilangan}
\begin{problem}
    (Afrika Selatan MO 2006) Sederhanakan pecahan $$\frac{2121212121210}{1121212121211}$$ ke bentuk paling sederhana.
\end{problem}

\begin{problem}
    (Afrika Selatan MO 2021) Tentukan bilangan asli terkecil dan terbesar dengan representasi digit yang berbentuk $ababa (a \neq 0)$ dan habis dibagi $11$.
\end{problem}

\begin{problem}
    (AIME 1984) Bilangan asli $n$ adalah bilangan asli terkecil yang habis dibagi oleh $15$ dan semua digitnya sama dengan $8$ atau $0$. Tentukan nilai dari $\frac{n}{15}$.
\end{problem}

\begin{problem}
    (Kualifikasi Chile MO 2017)
    Dengan menggunakan dua digit berbeda $2$ dan $d$, bilangan $90$ digit $m = 2d22d222d2222d...$ habis dibagi oleh $9$. Tentukan semua nilai yang mungkin dari $d$.
\end{problem}

\begin{problem}
    (El Salvador MO 2020) Perhatikan bilangan $$x_n =\overline{abaabaaab...b}\overline{\underbrace{a...a}_{n\,\, kali}{b}}$$ dengan $a\neq 0$ dan $b$ keduanya merupakan digit. Tentukan banyak bilangan asli $n$ dengan $1 \leq n \leq 2020$ sedemikian hingga $x_n$ habis dibagi oleh $11$ tak peduli apapun nilai $a$ dan $b$.
\end{problem}

\subsection{Beberapa Properti}
\begin{enumerate}
\item Untuk setiap bilangan bulat $a \neq 0$, $a|0$
\item Untuk setiap bilangan bulat $a \neq 0$, $a|a$
\item Untuk setiap bilangan bulat $a,b \neq 0$, jika $a|b$ dan $b|c$, maka $a|c$
\item Untuk setiap bilangan bulat $a \neq 0$, jika $a|b$ dan $a|c$, maka $a|\alpha b + \beta c$ untuk setiap bilangan bulat $\alpha, \beta$
\item Untuk setiap bilangan bulat $a,b \neq 0$ dengan $a|b$, maka $|a|\leq |b|$
\item Untuk setiap bilangan bulat $a,b \neq 0$ dengan $a|b$ dan $b|a$, maka $|a| = |b|$
\end{enumerate}
\subsection{Algoritma Pembagian}
Untuk setiap bilangan asli $a,b$, terdapat pasangan \textbf{unik} bilangan cacah $(q,r)$ yang memenuhi $$b = qa + r,  r<a$$ 
% \textbf{Bukti :} \\
% Jika $a > b$, maka $q = 0$ dan $r = b$\\
% Jika $a = b$, maka $q = 1$ dan $r = 0$\\
% Jika $a < b$, maka ambil $q$ sebagai bilangan bulat terkecil yang memenuhi $(q + 1)a > b$. Maka $qa \leq b$ dan $r = b - qa$ memenuhi $0 \leq r < a$.\\
% Untuk membuktikan keunikan, asumsikan terdapat dua pasangan $(q,r)$ dan $(q',r')$ yang memenuhi properti tersebut, maka $aq + r = aq' + r'$. Akibatnya, $a(q - q') = (r' - r)$ dan $a| r' - r$. Maka, $|r' - r| \geq a$ atau $|r' - r| = 0$. Kemungkinan pertama tidak mungkin dapat dicapai karena $0 \leq r,r' < a$. Maka $r' = r$ dan $q = q'$.

\subsection{Contoh soal}
\begin{problem}
    Carilah semua bilangan bulat $n$ yang memenuhi $n + 1 | 100$
\end{problem}

\begin{problem}
    Carilah semua bilangan bulat $h$ yang memenuhi $n + 1 | n^2 + 99$
\end{problem}

\subsection{Latihan Keterbagian}
\begin{problem}
    Carilah semua bilangan asli $n$ sehingga $n + 1| n^2 + 1$.
\end{problem}

\begin{problem}
    Carilah semua bilangan asli $n$ sehingga $n + 10| n^3 + 100$
\end{problem}

\begin{problem}
    (Finlandia MO 2001) Tentukan semua bilangan asli $n$ sedemikian hingga $n^2 + 2$ habis membagi $2 + 2001n$
\end{problem}

\begin{problem}
    (India PRMO 2014) Banyak bilangan asli $n$ di antara $1$ sampai $2014$ sehingga $\frac{8n}{9999 - n}$ merupakan bilangan bulat adalah ...
\end{problem}

\begin{problem}
    Diberikan sebuah barisan sepuluh bilangan asli naik tegas sedemikian hingga setiap bilangan membagi bilangan berikutnya. Jika diketahui jumlah kesepuluh bilangan tersebut adalah $1918$, maka semua kemungkinan bilangan kesepuluh dari barisan tersebut adalah?
\end{problem}

\begin{problem}
    Tentukan semua nilai bilangan asli $n$ sedemikian hingga $\lfloor \sqrt{n} \rfloor$ membagi $n$
\end{problem}

% \newpage

% \section{Bilangan Prima}
% Suatu bilangan disebut sebagai bilangan prima apabila hanya ada tepat satu bilangan $d>1$ yang membagi $p$, yakni $d = p$. Bilangan asli yang $n>1$ disebut sebagai komposit jika $n$ bukan bilangan prima.
% \subsection{Teorema Faktorisasi Prima}
% Setiap bilangan asli $n > 1$ dapat dipresentasikan secara unik sebagai perkalian dari bilangan- bilangan prima.\\
% Sebagai akibat dari teorema di atas, setiap biangan asli $n > 1$ dapat dinyatakan secara unik dalam bentuk $$n = p_1^{\alpha_1}p_2^{\alpha_2}...p_k^{\alpha_k}$$ dimana $p_1,p_2,...,p_k$ adalah bilangan prima berbeda dan $\alpha_1,\alpha_2,...,\alpha_k$ adalah bilangan asli.
% \subsection{Beberapa Properti yang Sering Digunakan}
% \begin{enumerate}
% \item Banyak bilangan prima adalah tak hingga
% \item Setiap bilangan prima yang lebih besar dari $3$ pasti berbentuk $6k \pm 1$ untuk suatu bilangan asli $k$
% \item Jika $a,b$ adalah bilangan bulat dan $p$ bilangan prima sehingga $p|ab$, maka $p|a$ atau $p|b$ (atau keduanya)
% \end{enumerate}

% \subsection{Latihan Bilangan Prima}
% \begin{problem}
%     Tentukan semua bilangan asli $n$ sedemikian hingga $3n - 4, 4n - 5$, dan $5n - 3$ semuanya merupakan bilangan prima
% \end{problem}

% \begin{problem}
%     (OSK 2011) Tentukan semua bilangan asli $p$ sedemikian hingga $p, p + 8$, dan $p + 16$ semuanya merupakan bilangan prima
% \end{problem}

% \begin{problem}
%     Tentukan semua bilangan asli $a, b$ sedemikian hingga $a^4 + 4b^4$ merupakann bilangan prima.
% \end{problem}

% \begin{problem}
%     (OSP 2010) Tentukan semua bilangan prima $p$ sehingga $p^2 + 7^3$ merupakan bilangan kubik.
% \end{problem}

% % \subsection{Soal-Soal}
% % \begin{enumerate}
% % \item Buktikan bahwa bilangan prima $> 3$ selalu berbentuk $6k \pm 1$ untuk sutau bilangan asli $k$.
% % \item Buktikan bahwa untuk $p>3$ prima, $24|p^2 - 1$
% % \item Tunjukkan bahwa untuk $a,b$ bilangan asli,  $a^4 + 4b^4$ merupakan bilangan prima hanya saat $a = b = 1$
% % \item Jika $p$ adalah bilangan prima ganjil dan $\dfrac{a}{b} = \dfrac{1}{1} + \dfrac{1}{2} + ... + \dfrac{1}{p - 1}$, buktikan bahwa $p|a$
% % \item Buktikan bahwa jika $a^n + 1$ dan $a$ bilangan asli yang lebih besar dari 1, maka $a$ genap dan $n = 2^k$ untuk suatu bilangan cacah $k$. Bilangan prima dalam bentuk $2^{2^k} + 1$ disebut sebagai bilangan prima Fermat.
% % \item Buktikan bahwa jika $a^n - 1$ dan $a$ bilangan asli yang lebih besar dari 1, maka $a = 2$ dan $n$ prima. Bilangan prima dalam bentuk $2^n - 1$ disebut sebagai bilangan prima Mersenne.
% % \item Misalkan $p,q,r$ adalah bilangan prima yang memenuhi $p^2 + q^2 = r^n$ dengan $n$ adalah bilangan asli. Tunjukkan bahwa $n = 1$
% % \end{enumerate}

% \section{FPB dan KPK (Bahasa inggris : GCD and LCM)}
% \subsection{FPB}
% Misalkan untuk bilangan asli $k$, himpunan $D_k$ adalah himpunan seluruh faktor positif dari $k$. $FPB(m,n)$ adalah anggota terbesar di himpunan $D_m \cap D_n$. Ketika $FPB(m,n) = 1$, kita menyebut $m$ dan $n$ relatif prima.
% \subsection{Properti FPB}
% \begin{enumerate}
% \item Jika $d = FPB(m,n), m = dm', n = dn'$, maka $FPB(m',n') = 1$
% \item Jika $d = FPB(m,n), m = d'm", n = d'n", FPB(m",n") = 1$, maka $d' = d$
% \item Jika $d'$ adalah faktor persekutuan dari $m$ dan $n$, maka $d'|d$
% \item Jika $m = p_1^{\alpha_1}p_2^{\alpha_2}...p_k^{\alpha_k}$ dan $n = p_1^{\beta_1}p_2^{\beta_2}...p_k^{\beta_k}$ dengan $\alpha_i, \beta_i \geq 0$ dan $\alpha_i + \beta_i \geq 1$ untuk $i =1,2,..., k$, maka $$FPB(m,n) = p_1^{min(\alpha_1, \beta_1)}p_2^{min(\alpha_2, \beta_2)}...p_k^{min(\alpha_k, \beta_k)}$$
% \item Jika $m = nq + r$, maka $FPB(m,n) = FPB(n,r)$
% \end{enumerate}
% \subsection{Algoritma Euclid}
% Merupakan sebuah algoritma yang berguna untuk mencari FPB dari dua bilangan dengan menggunakan Algoritma Pembagian secara berulang :
% \begin{align*}
%   m  &= nq_1 + r_1,  1\leq r_1 < n \\
%   n &= r_1q_2 + r_2, 1\leq r_2 < r_1 \\
%   r_1 &= r_2q_3 + r_3, 1\leq r_3 < r_2 \\
%   &... \\
%   r_{k - 2} &= r_{k - 1}q_k + r_k, 1\leq r_k < r_{k - 1} \\
%   r_{k - 1} &= r_kq_{k + 1} + r_{k + 1}, r_{k + 1} = 0
% \end{align*}
% Maka, $FPB(m,n) = FPB(n, r_1) = FPB(r_1, r_2) = ... = FPB(r_{k - 1}, r_k) = r_k$\\
% Sebagai akibat dari Algoritma Euclid ini, untuk setiap bilangan asli $m,n$ terdapat bilangan bulat $a,b$ sedemikian hingga $am + bn = FPB(m,n)$.
% \subsection{KPK}
% Misalkan untuk bilangan asli $k$, himpunan $M_k$ adalah himpunan seluruh kelipatan positif dari $k$. $KPK(m,n)$ adalah anggota terkecil di himpunan $M_m \cap M_n$. 
% \subsection{Properti KPK}
% \begin{enumerate}
% \item Jika $m = KPK(s,t), m = ss', m = tt'$, maka $FPB(s',t') = 1$
% \item Jika $m'$ adalah kelipatan persekutuan dari $s$ dan $t$ dan $m' = ss' = tt', FPB(s',t') = 1$, maka $m' = m$
% \item Jika $m'$ adalah kelipatan persekutuan dari $s$ dan $t$, maka $m|m'$
% \item Jika $m = p_1^{\alpha_1}p_2^{\alpha_2}...p_k^{\alpha_k}$ dan $n = p_1^{\beta_1}p_2^{\beta_2}...p_k^{\beta_k}$ dengan $\alpha_i, \beta_i \geq 0$ dan $\alpha_i + \beta_i \geq 1$ untuk $i =1,2,..., k$, maka $$FPB(m,n) = p_1^{max(\alpha_1, \beta_1)}p_2^{max(\alpha_2, \beta_2)}...p_k^{max(\alpha_k, \beta_k)}$$
% \item $FPB(m,n)\times KPK(m,n) = mn$
% \end{enumerate}
% \subsection{Latihan FPB dan KPK}
% \begin{problem}
%     Hitunglah $FPB(133, 198)$ dan $FPB(6436343, 12167)$
% \end{problem}

% \begin{problem}
%      Berapa banyak pasangan bilangan asli $(x,y)$ yang memenuhi $FPB(x,y) = 5!$ dan $KPK(x,y) = 50!$?    
% \end{problem}

% \begin{problem}
%     Definisikan $f(n) = FPB(n! + 1, (n + 1)!)$. Tentukan nilai dari $f(1) + f(2) + f(3) + ... + f(30)$
% \end{problem}

% \begin{problem}
%     (OSK 2021 Kemampuan Lanjut) Bilangan asli ganjil $b$ terbesar sedemikian hingga barisan bilangan asli $a_n = n^2 + 19n + b$ memenuhi \[FPB(a_n, a_{n + 1}) = FPB(a_{n + 2}, a_{n + 1})\] untuk setiap bilangan asli $n$ adalah ...
% \end{problem}

% \begin{enumerate}
% \item Hitunglah $FPB(133, 198)$ dan $FPB(6436343, 12167)$
% \item Misalkan $m$ dan $n$ adalah bilangan asli yang memenuhi $$FPB(m,n) + KPK(m,n) = m + n.$$ Buktikan bahwa salah satu dari kedua bilangan itu membagi bilangan yang lainnya.
% \item Carilah semua tripel bilangan asli $(a,b,c)$ sedemikian hingga $a^3 + b^3 + c^3$ habis dibagi $a^2b, b^2c,$ dan $c^2a$.
% \item Misalkan $x,y,z$ adalah bilangan asli yang memenuhi $$\dfrac{1}{x} - \dfrac{1}{y} = \dfrac{1}{z}$$ Misalkaan $h = FPB(x,y,z)$. Tunjukkan bahwa $hxyz$ dan $h(y - x)$ keduanya merupakan bilangan kuadrat.
% \item Berapa banyak pasangan bilangan asli $(x,y)$ yang memenuhi $FPB(x,y) = 5!$ dan $KPK(x,y) = 50!$?
% \item Diberikan bilangan asli $a,b$ yang memenuhi $$\dfrac{a + 1}{b} + \dfrac{b + 1}{a}$$ adalah bilangan asli. Tunjukkan bahwa $FPB(a,b) \leq \sqrt{a + b}$
% \end{enumerate}

% \section{Chinese Remainder Theorem}
% Misalkan $m_1,m_2,...,m_k$ adalah bilangan asli $\neq 1$ yang saling relatif prima. Maka, untuk setiap bilangan bulat $a_1,a_2,...,a_k$, sistem kekongruenan $x \equiv a_i$ (mod $m_i$), $i = 1,2,...,k$ memiliki solusi yang unik modulo $m = m_1m_2...m_k$, yakni $$\sum_{j = 1}^k \dfrac{m}{m_j}b_ja_j$$ dengan $b_j \equiv m_jm^{-1}$ (mod $m_j$)

\end{document}
